% GENERATED FILE. Edit canonical lessons, then run npm run generate:books.
% canonical-source-hash: fnv1a64:67c0fed7fcc16201
% canonical-lessons: GE-C252-ich-ging, GE-C252-ich-las, GE-C252-als, GE-C252-review

\chapter{The Told Story — ging, kam, las, als}
\label{ch:ge-the-told-story}
\hlchaptermodality{252}

\begin{chapteropening}
\textbf{By the end of this chapter:} \emph{I can tell what happened in the Präteritum with ich ging, ich kam and ich las, and place one event in the past with als.}

Complete GE-C252-review: two or three sentences of a past story in the Präteritum, one background with während and one event with als.
\end{chapteropening}

\section[ich ging, ich kam, ich sah]{ich ging, ich kam — the past a story is told in}

\label{lesson:GE-C252-ich-ging}



\begin{quote}
Say \textquotedblleft{}I said\textquotedblright{} in one word. (\emph{Ich sagte.}) Now \textquotedblleft{}I made\textquotedblright{}. (\emph{Ich machte.})

That \emph{-te} is the Präteritum, the past a German book or a fairy tale is told in. But the verbs a story leans on hardest, going, coming, seeing, do not take \emph{-te} at all.
\end{quote}

\begin{grammarlens}[title={Grammar Lens: the vowel changes, and nothing is added}]
Six verbs you already own, and their one-word past:

\begin{quote}
gehen — \textbf{ich ging}, I went kommen — \textbf{ich kam}, I came sehen — \textbf{ich sah}, I saw fahren — \textbf{ich fuhr}, I travelled, I drove nehmen — \textbf{ich nahm}, I took stehen — \textbf{ich stand}, I stood
\end{quote}

No \emph{-te} anywhere. The stem's vowel changes instead, and for \emph{ich} and for \emph{er} or \emph{sie} nothing is added: \emph{ich ging}, \emph{er ging}, just as \emph{ich sagte}, \emph{er sagte}. The plural adds \textbf{-en}: \emph{wir gingen}, \emph{sie sahen}. And \emph{du} adds \textbf{-st}: \emph{du gingst}.

These are the \textbf{strong} verbs. Which vowel each one takes is learned verb by verb, like English \emph{went}, \emph{came} and \emph{saw}.

A separable verb splits in the Präteritum as it does in the present: \emph{ankommen}, \textbf{ich kam in Berlin an}.
\end{grammarlens}

\begin{cousinweb}
Changing the vowel to make a past is older than German. The Germanic languages inherited it, and English kept it too: \emph{sing, sang}; \emph{come, came}. German calls it \textbf{Ablaut}, \textquotedblleft{}off-sound\textquotedblright{}.

\emph{Kommen} goes back to Old High German \emph{kweman} and \emph{sehen} to \emph{sehan}, which EtymDB 2.1 traces to the Proto-Germanic parents of English \emph{come} and \emph{see}. EtymDB does not list the past forms, so this is the usual account, given as one: Old High German already said \emph{quam}, \textquotedblleft{}came\textquotedblright{}, and \emph{sah}, \textquotedblleft{}saw\textquotedblright{}. \emph{Ging} is usually traced to a second old verb of going, \emph{gangan}.

The \emph{-te} of \emph{sagte} is the younger Germanic invention from chapter twenty-five, and new verbs still take it.
\end{cousinweb}

\subsection*{Your turn}
Say these aloud:

\begin{itemize}
\raggedright
  \item \emph{gehen}, \emph{ich ging}; \emph{kommen}, \emph{ich kam}
  \item \emph{sehen}, \emph{ich sah}; \emph{fahren}, \emph{ich fuhr}
  \item \emph{nehmen}, \emph{ich nahm}; \emph{stehen}, \emph{ich stand}
  \item \emph{ich ging}, \emph{er ging}, \emph{wir gingen}: nothing, nothing, then \emph{-en}
  \item \emph{Ich kam in Berlin an.} The prefix still goes last
  \item \emph{ich sagte}, then \emph{ich sah}: one with \emph{-te}, one with a new vowel
\end{itemize}

\begin{tcolorbox}[breakable,colback=teal!4,colframe=teal!35!black,title={Before you move on}]
\textquotedblleft{}I travelled\textquotedblright{}? (\emph{\textbf{Ich fuhr.}}) \textquotedblleft{}He took\textquotedblright{}? (\emph{\textbf{Er nahm}}, the same as \emph{ich nahm}.) \textquotedblleft{}We came\textquotedblright{}? (\emph{\textbf{Wir kamen.}}) What changes in a strong verb's past, if not the ending? (\textbf{The vowel.}) And what is the old name for that change? (\emph{\textbf{Ablaut.}})
\end{tcolorbox}
\section[ich las / ich habe gelesen]{ich las — the reading, and the having read}

\label{lesson:GE-C252-ich-las}



\begin{quote}
Say \textquotedblleft{}I went\textquotedblright{} and \textquotedblleft{}I saw\textquotedblright{}. (\emph{Ich ging. Ich sah.})

\emph{Lesen} is strong too. Its past is \textbf{ich las}.
\end{quote}

\begin{grammarlens}[title={Grammar Lens: one past, two English meanings}]
English has two pasts for one verb: \textquotedblleft{}I read\textquotedblright{}, finished, and \textquotedblleft{}I was reading\textquotedblright{}, still going on. German has no \textquotedblleft{}-ing\textquotedblright{} past. \textbf{Ich las} is both:

\begin{quote}
\emph{Ich las die Nachricht.} — I read the message. \emph{Ich las, und es regnete.} — I was reading, and it was raining.
\end{quote}

The sentence around the verb tells you which. A message gets read and is done. A rainy afternoon of reading is a scene: it was going on while something else happened.

That is how a German story sets its scene. \emph{Es war Samstag. Es regnete. Ich las.} Three Präteritum verbs, and none of them is finished yet. The weather verb takes the plain \emph{-te}, with one extra \emph{e} to make it sayable: \emph{es regnet}, \textbf{es regnete}. \emph{Warten} does the same: \emph{ich wartete}, I waited, I was waiting.
\end{grammarlens}

\begin{culture}
When German wants to say plainly that one thing was still going on, it reaches for \textbf{während}, \textquotedblleft{}while\textquotedblright{}. It works like \emph{weil}: the verb goes to the end of its clause.

\begin{quote}
\emph{Während ich las, klingelte das Telefon.} — While I was reading, the phone rang.
\end{quote}

The reading is the scene; the phone is the event that cut into it.

In conversation the same afternoon comes out in the Perfekt: \emph{ich habe gelesen}. The book and the story keep \emph{ich las}. You may also hear \emph{ich war am Lesen}, \textquotedblleft{}I was at the reading\textquotedblright{}, in everyday speech, especially in the west; it is colloquial, and a written story keeps the plain \emph{ich las}.
\end{culture}

\begin{cousinweb}
\textbf{lesen} $\leftarrow$ Old High German \emph{lesan} $\leftarrow$ Proto-Germanic \emph{*lesaną}, \textquotedblleft{}to gather\textquotedblright{} (EtymDB 2.1), and it still means \textquotedblleft{}to pick\textquotedblright{} in \emph{Weinlese}, the grape harvest. \emph{Las} is its past by the vowel change of the last lesson, already \emph{las} in Old High German by the usual account.
\end{cousinweb}

\subsection*{Your turn}
Say these aloud:

\begin{itemize}
\raggedright
  \item \emph{ich las}, then \textquotedblleft{}I read\textquotedblright{} and \textquotedblleft{}I was reading\textquotedblright{}: one German form
  \item a scene in three: \emph{Es war Samstag. Es regnete. Ich las.}
  \item \emph{Während ich las, klingelte das Telefon.}
  \item \emph{ich wartete}: the extra \emph{e} after the \emph{t}
  \item the conversation version: \emph{Ich habe gelesen.}
\end{itemize}

\begin{tcolorbox}[breakable,colback=teal!4,colframe=teal!35!black,title={Before you move on}]
How does German say \textquotedblleft{}I was reading\textquotedblright{} in a story? (\emph{\textbf{Ich las}}, the same as \textquotedblleft{}I read\textquotedblright{}.) Which word holds a scene open? (\emph{\textbf{Während.}}) And where does the verb go after it? (\textbf{To the end.}) \textquotedblleft{}It was raining\textquotedblright{}? (\emph{\textbf{Es regnete.}})
\end{tcolorbox}
\section[als, wenn]{als and wenn — once, or every time}

\label{lesson:GE-C252-als}



\begin{quote}
Say \textquotedblleft{}while I was reading\textquotedblright{}. (\emph{Während ich las.}) Then ask \textquotedblleft{}when are you coming?\textquotedblright{} (\emph{Wann kommst du?})

\emph{Wann} asks. Inside a story, German has two other words for \textquotedblleft{}when\textquotedblright{}, and it chooses by counting.
\end{quote}

\begin{grammarlens}[title={Grammar Lens: one occasion, or every occasion}]
For \textbf{one} occasion in the past, German says \textbf{als}:

\begin{quote}
\emph{Als ich in Berlin war, sah ich meine Großmutter.} — When I was in Berlin (that one time), I saw my grandmother.
\end{quote}

For something that happened \textbf{every time}, it says \textbf{wenn}, \textquotedblleft{}whenever\textquotedblright{}:

\begin{quote}
\emph{Wenn ich in Berlin war, besuchte ich meine Großmutter.} — Whenever I was in Berlin, I visited my grandmother.
\end{quote}

English says \textquotedblleft{}when\textquotedblright{} for both. Both German words work like \emph{weil}: the verb goes to the \textbf{end} of its clause, \emph{als ich in Berlin} \textbf{war}.
\end{grammarlens}

\begin{culture}
In the first sentence two verbs meet at the comma: \emph{… in Berlin} \textbf{war}, \textbf{sah} \emph{ich}. The whole \emph{als} clause is the first piece of the sentence, and the main verb comes second, so it follows the comma at once, with the subject behind it.

Turn the sentence round and the meeting disappears: \emph{Ich sah meine Großmutter, als ich in Berlin war.}

With the present, \emph{wenn} means \textquotedblleft{}when\textquotedblright{} or \textquotedblleft{}whenever\textquotedblright{}: \emph{Wenn ich in Berlin bin, besuche ich sie.} \emph{Als} belongs to one occasion that is over.
\end{culture}

\begin{cousinweb}
\textbf{wenn} and \textbf{wann} were one word in Old High German, \emph{wanne} or \emph{wenne}; EtymDB 2.1 traces \emph{wann} to Proto-Germanic \emph{*hwannē}, the parent of English \emph{when}. It lists no parent for \emph{wenn}, so the split of one old word into a question word and a joining word is given as the usual account.

\textbf{als} $\leftarrow$ Old High German \emph{alsō} (EtymDB 2.1), worn down to one syllable. EtymDB gives no meaning for it; by the usual account it was \emph{al} + \emph{sō}, \textquotedblleft{}all so, just so\textquotedblright{}.
\end{cousinweb}

\subsection*{Your turn}
Say these aloud:

\begin{itemize}
\raggedright
  \item \emph{Als ich in Berlin war, sah ich meine Großmutter.} One time
  \item \emph{Wenn ich in Berlin war, besuchte ich meine Großmutter.} Every time
  \item the two verbs at the comma: \emph{war, sah}
  \item turned round: \emph{Ich sah meine Großmutter, als ich in Berlin war.}
  \item \emph{Als ich las, klingelte das Telefon.} One time
\end{itemize}

\begin{tcolorbox}[breakable,colback=teal!4,colframe=teal!35!black,title={Before you move on}]
One time in the past: which word? (\emph{\textbf{Als.}}) Every time? (\emph{\textbf{Wenn.}}) A question about time? (\emph{\textbf{Wann.}}) Where does the verb go after \emph{als}? (\textbf{To the end.}) And what comes straight after the comma when the \emph{als} clause goes first? (\textbf{The main verb.})
\end{tcolorbox}
\section[(review)]{Review — the scene, and what cut into it}

\label{lesson:GE-C252-review}



\begin{quote}
Say \textquotedblleft{}I went\textquotedblright{}, \textquotedblleft{}I came\textquotedblright{} and \textquotedblleft{}I saw\textquotedblright{}. (\emph{Ich ging. Ich kam. Ich sah.}) Then \textquotedblleft{}I waited\textquotedblright{}. (\emph{Ich wartete.})
\end{quote}

\begin{grammarlens}[title={What you've built}]
Three lessons, and the past of a story is in your hands. Two ways to make it: the \emph{-te} of \emph{sagte} and \emph{wartete}, and the vowel change of \emph{ging}, \emph{kam}, \emph{sah}, \emph{fuhr}, \emph{nahm}, \emph{stand} and \emph{las}. One form for \textquotedblleft{}I read\textquotedblright{} and \textquotedblleft{}I was reading\textquotedblright{}, with the sentence around it deciding which. And \textbf{während} to hold a scene open while something happens.

Then the moment itself: \textbf{als} for one occasion, \textbf{wenn} for every time.

\begin{quote}
\emph{Während ich las, klingelte das Telefon.} \emph{Als ich in Berlin ankam, regnete es.}
\end{quote}

The second one is worth hearing twice. \emph{Ankommen} splits in a main clause: \emph{ich kam an}. After \emph{als}, the verb goes to the end, and the prefix joins it again: \emph{als ich in Berlin} \textbf{ankam}. Then verb, comma, verb: \emph{ankam, regnete}.
\end{grammarlens}

\subsection*{Your turn}
Say these aloud:

\begin{itemize}
\raggedright
  \item a scene in three: \emph{Es war Samstag. Es regnete. Ich las.}
  \item \emph{Während ich las, klingelte das Telefon.}
  \item \emph{Ich kam in Berlin an}, then \emph{Als ich in Berlin ankam, regnete es.}
  \item every time: \emph{Wenn ich in Berlin war, besuchte ich meine Großmutter.}
  \item one of your own: a scene with \emph{während}, and one thing that happened, with \emph{als}
\end{itemize}

\begin{tcolorbox}[breakable,colback=teal!4,colframe=teal!35!black,title={Before you move on}]
\textquotedblleft{}I took\textquotedblright{}? (\emph{\textbf{Ich nahm.}}) \textquotedblleft{}I was reading\textquotedblright{}? (\emph{\textbf{Ich las.}}) One occasion in the past: \emph{als} or \emph{wenn}? (\emph{\textbf{Als.}}) What happens to the \emph{an} of \emph{ankommen} after \emph{als}? (\textbf{It joins the verb again, at the end:} \emph{\textbf{ankam}}.)
\end{tcolorbox}
